\begin{ex}
Tìm hai số có tổng là 70 và hiệu là 10. Số lớn là bao nhiêu?
\shortans{40}
\loigiai{
Số lớn = $(70 + 10) : 2 = 40$.
}
\end{ex}

\begin{ex}
Tìm hai số có tổng là 70 và hiệu là 10. Số bé là bao nhiêu?
\shortans{30}
\loigiai{
Số bé = $(70 - 10) : 2 = 30$.
}
\end{ex}

\begin{ex}
Tổng của hai số là 100, hiệu của hai số là 20. Số lớn là bao nhiêu?
\shortans{60}
\loigiai{
$(100 + 20) : 2 = 60$.
}
\end{ex}

\begin{ex}
Tổng của hai số là 100, hiệu của hai số là 20. Số bé là bao nhiêu?
\shortans{40}
\loigiai{
$(100 - 20) : 2 = 40$.
}
\end{ex}

\begin{ex}
Hai bạn Lan và Mai có tất cả 48 cái nhãn vở, Lan có nhiều hơn Mai 8 cái. Hỏi Lan có bao nhiêu cái nhãn vở?
\shortans{28}
\loigiai{
$(48 + 8) : 2 = 28$ cái.
}
\end{ex}

\begin{ex}
Hai bạn Lan và Mai có tất cả 48 cái nhãn vở, Lan có nhiều hơn Mai 8 cái. Hỏi Mai có bao nhiêu cái nhãn vở?
\shortans{20}
\loigiai{
$(48 - 8) : 2 = 20$ cái.
}
\end{ex}

\begin{ex}
Lớp 4A có 36 học sinh, số học sinh nữ nhiều hơn số học sinh nam là 4 bạn. Hỏi lớp 4A có bao nhiêu học sinh nữ?
\shortans{20}
\loigiai{
$(36 + 4) : 2 = 20$ bạn nữ.
}
\end{ex}

\begin{ex}
Lớp 4A có 36 học sinh, số học sinh nữ nhiều hơn số học sinh nam là 4 bạn. Hỏi lớp 4A có bao nhiêu học sinh nam?
\shortans{16}
\loigiai{
$(36 - 4) : 2 = 16$ bạn nam.
}
\end{ex}

\begin{ex}
Một mảnh đất hình chữ nhật có nửa chu vi là 40 m, chiều dài hơn chiều rộng 10 m. Chiều dài mảnh đất là bao nhiêu mét?
\shortans{25}
\loigiai{
$(40 + 10) : 2 = 25$ m.
}
\end{ex}

\begin{ex}
Một mảnh đất hình chữ nhật có nửa chu vi là 40 m, chiều dài hơn chiều rộng 10 m. Chiều rộng mảnh đất là bao nhiêu mét?
\shortans{15}
\loigiai{
$(40 - 10) : 2 = 15$ m.
}
\end{ex}

\begin{ex}
Hai thùng chứa tất cả 120 lít dầu. Thùng thứ nhất chứa nhiều hơn thùng thứ hai 20 lít dầu. Thùng thứ nhất chứa bao nhiêu lít dầu?
\shortans{70}
\loigiai{
$(120 + 20) : 2 = 70$ lít.
}
\end{ex}

\begin{ex}
Hai thùng chứa tất cả 120 lít dầu. Thùng thứ nhất chứa nhiều hơn thùng thứ hai 20 lít dầu. Thùng thứ hai chứa bao nhiêu lít dầu?
\shortans{50}
\loigiai{
$(120 - 20) : 2 = 50$ lít.
}
\end{ex}

\begin{ex}
Tuổi bố và tuổi con cộng lại được 48 tuổi, bố hơn con 28 tuổi. Hỏi bố bao nhiêu tuổi?
\shortans{38}
\loigiai{
$(48 + 28) : 2 = 38$ tuổi.
}
\end{ex}

\begin{ex}
Tuổi bố và tuổi con cộng lại được 48 tuổi, bố hơn con 28 tuổi. Hỏi con bao nhiêu tuổi?
\shortans{10}
\loigiai{
$(48 - 28) : 2 = 10$ tuổi.
}
\end{ex}

\begin{ex}
Hai kho thóc có tất cả 500 tấn thóc. Kho A có ít hơn kho B là 60 tấn. Kho B có bao nhiêu tấn thóc?
\shortans{280}
\loigiai{
$(500 + 60) : 2 = 280$ tấn.
}
\end{ex}

\begin{ex}
Hai kho thóc có tất cả 500 tấn thóc. Kho A có ít hơn kho B là 60 tấn. Kho A có bao nhiêu tấn thóc?
\shortans{220}
\loigiai{
$(500 - 60) : 2 = 220$ tấn.
}
\end{ex}

\begin{ex}
Một hình chữ nhật có chu vi là 60 cm, chiều dài hơn chiều rộng 6 cm. Chiều dài hình chữ nhật đó là bao nhiêu xăng-ti-mét?
\shortans{18}
\loigiai{
Nửa chu vi = 30 cm. Dài = $(30 + 6) : 2 = 18$ cm.
}
\end{ex}

\begin{ex}
Một hình chữ nhật có chu vi là 60 cm, chiều dài hơn chiều rộng 6 cm. Chiều rộng hình chữ nhật đó là bao nhiêu xăng-ti-mét?
\shortans{12}
\loigiai{
Rộng = $(30 - 6) : 2 = 12$ cm.
}
\end{ex}

\begin{ex}
Hai số có tổng bằng số lớn nhất có hai chữ số (99) và hiệu bằng số lẻ bé nhất có hai chữ số (11). Số lớn là bao nhiêu?
\shortans{55}
\loigiai{
$(99 + 11) : 2 = 55$.
}
\end{ex}

\begin{ex}
Hai số có tổng bằng 99 và hiệu bằng 11. Số bé là bao nhiêu?
\shortans{44}
\loigiai{
$(99 - 11) : 2 = 44$.
}
\end{ex}

\begin{ex}
Tổng của hai số là 84, nếu lấy số lớn trừ đi số bé thì được 14. Số lớn là bao nhiêu?
\shortans{49}
\loigiai{
$(84 + 14) : 2 = 49$.
}
\end{ex}

\begin{ex}
Tổng của hai số là 84, nếu lấy số lớn trừ đi số bé thì được 14. Số bé là bao nhiêu?
\shortans{35}
\loigiai{
$(84 - 14) : 2 = 35$.
}
\end{ex}

\begin{ex}
Hai lớp 4A và 4B thu gom được tất cả 250 kg giấy vụn. Lớp 4A thu gom được nhiều hơn lớp 4B là 30 kg. Lớp 4A thu gom được bao nhiêu ki-lô-gam giấy vụn?
\shortans{140}
\loigiai{
$(250 + 30) : 2 = 140$ kg.
}
\end{ex}

\begin{ex}
Lớp 4B thu gom được bao nhiêu ki-lô-gam giấy vụn trong bài toán trên?
\shortans{110}
\loigiai{
$(250 - 30) : 2 = 110$ kg.
}
\end{ex}

\begin{ex}
Hai số có tổng bằng 1000 và hiệu bằng 200. Tích của hai số đó có tận cùng bằng bao nhiêu chữ số 0?
\shortans{4}
\loigiai{
Hai số là 600 và 400. Tích = $600 \times 400 = 240000$ (có 4 chữ số 0).
}
\end{ex}
